\documentclass[10pt]{article}
\usepackage{amsmath,amssymb,amsthm}
\newtheorem{theorem}{Theorem}
\title{ARC-SELECT-INDEX252: Virtual Child Input Addressing Audit}
\author{LSPV-4 Resource Auditor}
\date{30 October 2026}
\begin{document}\maketitle
\noindent\textbf{Status: PROVED/EXACT restricted directory; KNOWN FLMV lower bound; COMPUTATIONAL finite control; universal ARC and P versus NP OPEN.} Independent review pending.

\section*{Target and four falsifiers}
The single integrated task is a universal, fully charged SAT-to-SAT arbitrary-decider transformation; no such transformation is constructed. Tests: certify the exact FLMV infinite envelope; audit checkpointed virtual child input access; quantify repeated per-child full scans; identify an easy structured escape. Prior LSPV-4 FLMV-CHILD-GATE251 established the additive-versus-multiplicative guest-cost recurrence and the strict22 plain-cofactor size barrier. Neither result is claimed as new.

\section*{Primary FLMV theorem and strict entry}
Let $f_0=1$, $f_{k+1}=cf_k/(1+f_k)$ and
$d_k(c)=\min\{(1/f_k-c)/k,c^{-k}\}$ for $k\ge1$. The original printed Theorem 5.7 uses $k\ge1$. Since $d_j(c)\le c^{-j}$, if $b=\max_{i\le K}d_i(c)>c^{-(K+1)}$, all later indices are strictly inferior. Exact rational checking certifies $d_*(6/5)=23/45$ and
$d_*(241/200)=11642479/23232400$, both at $k=2$, with tail certified after $K=3$. A contradiction requires one matched uniform deterministic total indexed-input SAT machine with exponents $K<241/200$ and $L<11642479/23232400$, not two independently optimized machines. At the chosen point the writable-space exponent margin above $1/2$ is $26279/23232400$.

\section*{New restricted physical child-address theorem}
Consider an explicit source-order cofactor emitter that deletes satisfied source clauses. Let $b\in\{0,1\}^m$ mark retained original clause slots. Child clause $j$ comes from original slot $\operatorname{select}_b(j)$. Grant read-only access to the deletion bitmap, charging each accessed bit as one source query. A \emph{single-level block-prefix directory} divides the bitmap into $p=\lceil m/h\rceil$ blocks of length at most $h$ and stores only one full-width cumulative rank at each block end; it stores no miniblock summary or other index. Set $w=\lceil\log_2(m+1)\rceil$. The directory uses exactly $S=pw$ indexed bits, after one complete scan of $m$ marks. One query binary-searches the block counters and scans at most $h$ bitmap bits.

\begin{theorem}[Restricted checkpoint resource product]
For the declared single-level prefix-directory architecture, some legal source bitmap and select query require at least $h-1$ within-block bitmap probes, even when the block's total number of survivors is known to be one. Thus its worst-case query cost $Q$ satisfies $S(Q+1)\ge mw$.
If original fixed-slot source length is $N=\Theta(m\log m)$, realizing $S=O(N^{s+o(1)})$ and $Q=O(N^{e+o(1)})$ requires $s+e\ge1$. In particular with proposed root directory space exponent $s=1/2$, a bad query requires $N^{1/2-o(1)}$ source bitmap probes. This restricted architecture cannot simultaneously meet the hypothetical ARC $s=1/2$ and per-guest-input access $e<1/400$ budgets.
\end{theorem}
\begin{proof}
Consider a block with exactly one retained bit. Its sole one-position is not revealed by the scalar prefix counters. An adaptive exact bit-probe procedure can receive zero for its first $h-1$ distinct inspected positions, requiring at least $h-1$ probes in the worst case. Since $S=\lceil m/h\rceil w$, multiplication gives $S(Q+1)\ge mw$. Substituting the source encoding and taking exponents gives $s+e\ge1$.
\end{proof}
This is an index-architecture theorem, \emph{not} a general rank/select or arbitrary-SAT lower bound. On the explicit periodic bitmask $b_i=1$ iff $i\bmod4\ne3$, child ordinal $j$ maps to zero-based source slot $j-1+\lfloor(j-1)/3\rfloor$ without bitmap probes. Sophisticated succinct rank/select or source-aware emitter numbering may also escape. A fully materialized bitmap in writable memory costs $m$ bits; Python finite tests do materialize it and are not native low-space FLMV executions. If $N^{3/5}$ children each independently rebuild a length-$m=\Theta(N/\log N)$ deletion mask, this \emph{particular full-rescan implementation} expends $N^{8/5-o(1)}$ source reads. Incremental update and sharing are possible alternatives.

\section*{Both-label finite experiment}
The source-indexed Python program verified W1 strict signed exact22 UNSAT with zero models, its matched two-sign-swapped strict22 SAT neighbor with 168 models, H1 SAT with four, XOR3 SAT with four, contradictory XOR3 UNSAT with zero, and Horn SAT with one. Original global 0/1 assignments, complement and signed occurrence targets remain adversarial controls. Each original cofactor under pins $(x_1=0)$, $(x_1=1)$ and $(x_1=0,x_2=1)$ was compared against physically materialized original-order source clause records using checkpoint lengths $1,2,3,4,8,16,32$. Independent Node code passed 126 directory configurations and 987 retained-clause select queries, plus exact FLMV arithmetic. Six sparse lengths 64 through 2048 reproduced worst-block scans under square-root directory budgets; periodic sources avoided them.
\[
\begin{array}{lrrrr}
\text{Source}&\text{models}&\text{base DPLL nodes}&\text{two child nodes}&\text{base/child literal reads}\\
\hline
\text{W1 UNSAT}&0&151&150&14045/12576\\
\text{W1 signed-swap SAT}&168&19&26&1473/1782\\
\text{H1 SAT}&4&7&9&184/186\\
\text{XOR SAT}&4&3&4&34/20\\
\text{XOR UNSAT}&0&7&6&160/62
\end{array}
\]
Child solver reads exclude two required source-to-child conversion scans, output writes, rank/select indexing, all policy and failed-trial work, and certificate checking. These are diagnostic logical operations, not native FLMV steps or a claimed net improvement. One-pin W1 children occupy 1600 or 1608 bits in this literal plain-text encoding versus 1808 source bits: no fixed-power shrink. Same-source SAT and UNSAT outcomes were independently checked by exhaustive enumeration.

\section*{Sources and status}
This run read active governance UNIVERSAL-PARAMETERIZED-FLMV-REDIRECT1, universal contractor, AFFINE-SMALLBIAS-DESCENT-GATE2, CONSTRUCTION-STRENGTHEN17, ARC-LINEAR22-SATURATION1, LSPV-4 CHILD-GATE251, peer LSPV-2 COFACTOR-VOLUME-ARC252, current root, AGENTS, CODEX, research skill, coordination issue 3, original FLMV 2005 and primary rank/select work. The affine obstruction remains a restricted semantic result, not universal SAT hardness. Full source-indexed Python and raw data are retained in the reproducibility bundle; its SHA-256 raw hash is \texttt{c50bae1e9531093e48bc80724154641493168bb80b2307d955f5ef9f7c689786}. The next OPEN task is an explicit arbitrary-bit, physically smaller SAT-child emitter with paid address translation and two-level hereditary correctness.

\begin{thebibliography}{9}
\bibitem{flmv} Fortnow, Lipton, van Melkebeek, Viglas, Time-Space Lower Bounds for Satisfiability, JACM 52(6), 2005.
\bibitem{kurpicz} F. Kurpicz, Engineering Compact Data Structures for Rank and Select Queries on Bit Vectors, arXiv:2206.01149.
\bibitem{dos} Dreier, Ordyniak, Szeider, SAT Backdoors: Depth Beats Size, JCSS 142 (2024).
\end{thebibliography}
\end{document}
% END_OF_VERIFIED_REPORT::LSPV-4::ARC-SELECT-INDEX252
